\documentclass[a4paper,fleqn]{article}
%در صورت نیاز به فراخوانی بسته‌های اضافی، آن‌ها را در فایل commands.tex و قبل از
%بسته زی‌پرشین فراخوانی کنید.
\input{commands}
\begin{document}
% دستوری جهت حذف کردن شماره صفحه و سربرگ، در صورت وجود (فقط در صفحه جاری)
\thispagestyle{empty}
\begin{center}
	{\Naz
	باسمه تعالی
	\\[.1cm]
	\includegraphics[height=3cm]{iut_logo.png}
	%دستوری برای کم کردن فاصله بین لوگو و خط پایین آن
	\vspace{-2mm}
	
	{\large
	دانشگاه صنعتی اصفهان\\
	}
	{\small
	دانشکده برق و کامپیوتر
	}
	%دستوری برای تعیین فاصله بین دو خط
	\\[.3cm]
	{\large
		\textbf{پیشنهاد موضوع تحقیقاتی برای پایان‌نامه کارشناسی ارشد}	
		\\[.2cm]
	}

%%{
%%  \fontsize{9}{10.8}\selectfont دانشجوی محترم: لطفاً 
%%	\textbf{شیوه نامه تدوین گزارش پیشنهاد موضوع تحقیق پایان‌نامه کارشناسی ارشد}	را از وب سایت دانشکده دريافت نمایید.	نسخه 94/06/31
%%}

}

\end{center}
%%%%%%%%%%%%%%%%%%%%%%%%%%%%%%%%%%%%%%%%%%%
%\setlength{\parindent}{0pt}
%\vspace{-0.25cm}
%%%%%%%%%%%%%%مشخصات فردی را وارد کنید
{
\titr
\textbf{مشخصات دانشجو:}
}
{\Naz
\begin{center}
	\scalebox{1.1}{
		\begin{tabular}{|r|r|r|r|}
			
			\hline
			نام و نام خانوادگی: نوید نصیری&
			شماره دانشجويي: 40214394&
			رشته: مهندسی کامپیوتر&
			گرایش: نرم‌افزار
			\\ \hline
			\multicolumn{2}{|r|}{گرایش/رشته  دوره کارشناسی: مهندسی کامپیوتر-سیستم‌های هوشمند}&
			\multicolumn{2}{|r|}{شماره تلفن همراه: 09215312706}
			\\	\hline
			
		\end{tabular}}
\end{center}}
%	\vfill
%%%%%%%%%%%%%%عنوان پروپوزال را وارد کنید
{
\titr
\textbf{مشخصات تحقیق پیشنهادی:}
}
{\Naz
\begin{center}
	\scalebox{0.97}{
    \begin{tabular}{|c|c|}
	\hline
	عنوان فارسی & دوقلو‌های مجازی چندعاملی
	\\ \hline  
	عنوان انگلیسی & \lr{Multi Agent Digital Twin}
	\\ \hline      
	\end{tabular}}
\end{center}
}
%%%%%%%%%%%%%% کلمات کلیدی گزارش
%کلمات کلیدی خود را وارد کنید.
\begin{center}
	\scalebox{1.}{
		\begin{tabular}{|r|r|r|r|}
			\hline
			کلمات کلیدی گزارش:&1- دوقلو‌های مجازی&2- سیستم‌ چند‌عاملی &3- هوش مصنوعی
			\\ \hline  
		\multicolumn{1}{|l|}{\lr{3- AI}} &
		\multicolumn{1}{|l|}{\lr{2- Multi Agent System}} &
		\multicolumn{1}{|l|}{\lr{1- Digital Twins}} &
		\multicolumn{1}{|l|}{Keywords:}
			\\ \hline
	\end{tabular}}
\end{center}
%%%%%%%%%%%%%%مشخصات استاد راهنما و مشاور را وارد کنید، موارد وارد شده برای مثال است.
{\titr
	\textbf{اساتید راهنما و مشاور:}}
\begin{center}
	\scalebox{0.89}{
		\begin{tabular}{|r|r|r|r|r|r|}
			\hline
			\multicolumn{1}{|c|}{سمت} &
			\multicolumn{1}{c|}{نام و نام خانوادگی} &
			\multicolumn{1}{c|}{درصد مشارکت} &
			\multicolumn{1}{c|}{دانشگاه/مؤسسه} &
			\multicolumn{1}{c|}{مرتبه علمی} &
			\multicolumn{1}{c|}{تخصص مرتبط با پایان نامه}
			\\ \hline
			\multicolumn{1}{|r|}{استاد راهنما} &
			دکتر  مسعودرضا هاشمی&۱۰۰ درصد&دانشگاه صنعتی اصفهان&استاد& شبکه‌های کامپیوتری
			\\ \hline
			\multicolumn{1}{|r|}{*استاد مشاور} &
			دکتر  عارف کریمی‌ افشار&۱۰۰ درصد&دانشگاه صنعتی اصفهان&استادیار& اینترنت اشیاء
			\\ \hline
			\multicolumn{1}{|r|}{*استاد راهنمای دوم} &
			&&&&
			\\ \hline
			\multicolumn{1}{|r|}{*استاد مشاور دوم} &
			&&&&
			\\ \hline
	\end{tabular}}
\end{center}
%	\vfill

\newpage
\setlength{\parindent}{0pt}
%\pagebreak
% دستوری جهت حذف کردن شماره صفحه و سربرگ، در صورت وجود (فقط در صفحه جاری)
\thispagestyle{empty}

\clearpage

	%%%%%%%%%%%سربرگ گزارش
	% دستوری جهت حذف کردن شماره صفحه و سربرگ، در صورت وجود (فقط در صفحه جاری)
\thispagestyle{empty}
\newgeometry{left=3cm,right=3cm,top=2cm}
{\Naz
\begin{center}
	باسمه تعالی
	\\[.1cm]
	\includegraphics[height=3cm]{iut_logo.png}
	%دستوری برای کم کردن فاصله بین لوگو و خط پایین آن
%	\vspace{-2mm}
	\vspace{0.2cm}
	
	{\large
		دانشگاه صنعتی اصفهان\\
	}
	{\small
		دانشکده برق و کامپیوتر
	}
	%دستوری برای تعیین فاصله بین دو خط
	\\[1.cm]
	{\LARGE
		\textbf{گزارش پيشنهاد موضوع تحقيق پایان‌نامه کارشناسی ارشد}	
		\\[1.3cm]
	}
%	در اینجا نسخه اول یا نهایی بودن گزارش خود را معین نمایید.
	{\Large
		\textbf{نسخه اول }
		\\[1.6cm]
	}
%	رشته خود را مشخص نمایید.
	{\LARGE
		\textbf{مهندسی کامپیوتر/نرم‌افزار}
		\\[1.2cm]
	}
%	عنوان را وارد کنید.
	{\titr
		{\LARGE
			\textbf{دوقلو‌های مجازی چندعاملی}
		\\[2.5cm]
		}
	}
%	نام خود را بنویسید
	{\LARGE
		\textbf{نوید نصیری}
		\\[2cm]
	}
%		نام استاد(های) راهنمای خود را وارد کنید. موارد نوشته شده برای مثال اند
	{\Large
		دکتر مسعودرضا هاشمی
%		\\
%		نام استاد مشاور را وارد کنید.
		\\[3.cm]
	}
	{\Large
	\textbf{تاریخ}
	\\[.4cm]
	\textbf{شهریور ۱۴۰۳}
	}
\end{center}}
\clearpage
% دستوری جهت حذف کردن شماره صفحه و سربرگ، در صورت وجود (فقط در صفحه جاری)
\thispagestyle{empty}
		

%*****************************************************************
%% تنظیم مناسب صفحه و فونت برای متن اصلی پایان‌نامه
\newgeometry{left=2.5cm,right=3cm,top=3cm,bottom=2.5cm,includehead=false,headsep=1cm,footnotesep=.5cm}
\settextfont{XB Zar}\fontsize{12}{6}\selectfont
\setlatintextfont{Times New Roman}\fontsize{11}{6}
\baselineskip=.9cm

\pagenumbering{arabic}
% Moving page number to top right
\pagestyle{myheadings}

\newpage

\section{موضوع کلی و زمینه اصلی تحقیق}

در دنیای امروز، با پیشرفت تکنولوژی، هزینه نگهداری و استفاده از آن نیز روز به روز افزایش می‌یابد. این افزایش هزینه مارا به بررسی روش‌ها و حالت‌های مختلف در دنیایی غیر از دنیای واقعی وا می‌دارد. ایجاد یک محیط مجازی از موجودیت‌های واقعی و فیزیکی یا به عبارتی دوقلوی مجازی\LTRfootnote{Digital Twin} 
می‌تواند این کاهش هزینه‌ را با خود به همراه داشته باشد. امروزه، مجازی‌سازی اشیاء، موجودیت‌ها و فرایند‌ها، این کمک را می‌کند که بتوان بدون نیاز به ایجاد تغییر و دستکاری در واقعیت، با مشاهده و بررسی دوقلوی مجازی آن، بهترین تصمیم گرفته و اجرا شود. 
با این وجود، یک حقیقت غیرقابل انکار این است که یک فرایند واقعی، از ارتباط و تاثیر موجودیت‌ها، اشیاء و فرایند‌های کوچک‌تر بر یکدیگر ایجاد می‌شود و نمی‌توان این مورد را نادیده گرفت. لذا مطالعه بر نوع ارتباط چند عامل\LTRfootnote{Multi Agent}
، مدیریت،‌ تبادل داده و تصمیم‌گیری درست بر اساس داده‌ها و هوش مصنوعی، می‌تواند به دوقلو‌های مجازی در دنیای مجازی کمک کند تا تاثیر بیشتر در واقعیت داشته‌ باشند.
پس تکنولوژی دوقلو‌های مجازی با استفاده از سیستم‌های چندعاملی، می‌تواند توسعه یافته و قابلیت و کارایی بالاتری را در واقعیت فراهم سازد.


\section{مرور کارها و نتایج  قبلی }

بر اساس مقاله 
\cite{b1}
دوقلو‌های مجازی معمولا از سه جزو ایجاد می‌شود. فضای واقعی دارای شئ واقعی، فضای مجازی دارای شئ مجازی و ارتباط میان شئ مجازی و واقعی و انتقال داده ‌میان آن‌ها.

طبق ادعایی که در مقاله‌های 
\cite{b2}
و
\cite{b3}
صورت گرفته‌است، دوقلو‌های مجازی تنها یک شبیه‌سازی ساده نیست. بلکه یک مدل زنده، هوشمند و در حال تکامل است. دوقلوی مجازی، چرخه عمر همتای حقیقی خود را دنبال می‌کند تا بتواند آن را نظارت\LTRfootnote{Monitor}
و کنترل\LTRfootnote{Control}
کرده و همچنین کارکرد آن را بهبود بخشد. دوقلوی مجازی به صورت مستمر با فضا و محیط واقعی در ارتباط بوده و به صورت بی‌درنگ با آن، داده‌ بین یکدیگر ردوبدل می‌کند.


 مقاله 
 \cite{b4}
 به بررسی دوقلو‌های مجازی در اینترنت وسایل‌ نقلیه\LTRfootnote{Internet of Vehicles (IoF)}
 پرداخته است. در واقع این مقاله، یک سکو\LTRfootnote{Platform}
 جهت به اشتراک گذاری داده، بین وسایل نقلیه موجود در سطح شهر را معرفی می‌کند. در این مقاله، یک دوقلوی دیجیتال، به صورت یک بازتاب از وسیله نقلیه در فضای ابری تعریف می‌شود که داده‌ها را با وسیله نقلیه در هر لحظه همگام می‌کند. داشتن این دوقلوی مجازی در بستر ابری، باعث می‌شود که بتوان با به اشتراک‌گذاری داده‌ها بدون محدودیت فاصله و پهنای باند، کارایی رانندگی را افزایش داده و تصمیمات دقیق‌تری بر اساس داده‌ها گرفته‌ شود.\\
در نهایت مقاله 
\cite{b5}
یک چارچوب\LTRfootnote{Framework}
 به نام 
\lr{MADTwin}\LTRfootnote{Multi Agent Digital Twin}
ارائه می‌دهد که برای مدل‌سازی دوقلوی دیجیتال برپایه عامل‌\LTRfootnote{Agent}
مناسب است. از این چارچوب برای هوشمند‌سازی یک انبار‌ توسط ربات‌ها استفاده کرده تا بتوان به طور بهینه هر ربات را به سوی اهداف محلی خود و اهداف کل سیستم به صورت مشارکتی یا مستقل هدایت کند.


\section{ضرورت انجام، روش و اهداف تحقیق پیشنهادی}
داشتن یک دوقلوی مجازی از یک یا چند فرایند، پدیده یا اشیاء، امکان توسعه و طراحی محصول جدید را تسهیل می‌کند. در واقع با کمک این فناوری، طراحان و مهندسان می‌توانند ایده‌های نوآورانه خود را پیش از پیاده‌سازی واقعی تست و ارزیابی کرده و سپس به سراغ پیاده‌سازی بروند. این امر به کاهش خطرات و هزینه‌های توسعه کمک چشم‌گیری می‌کند.

دوقلوی مجازی به وسیله توسعه الگوریتم‌های هوشمند، می‌تواند نمایی دقیق و کارآمد از سیستم‌های فیزیکی را در محیط‌های دیجیتال ارائه داده و تصمیم‌گیری‌های مفید و بهینه‌ای را انجام دهد. همچنین چندین دوقلوی مجازی هوشمند، می‌توانند با تعامل و همکاری باهم، تصمیمات گروهی مناسبی را اتخاذ کرده و در جهت بهبود فرایند‌ اصلی گام نهند.

از اهداف اصلی دوقلوی مجازی چندعاملی، می‌توان به افزایش کارایی و کاهش پیچیدگی مدیریت یک سامانه اشاره کرد. با تعامل چندین دوقلوی مجازی ساده‌تر و در نهایت ایجاد یک دوقلوی مجازی یکپارچه، مدیریت ساده‌تری برای کل سامانه فراهم شده به طوری که کل آن قابل خودکارسازی نیز باشد. همچنین این سادگی در مدیریت باعث کاهش قابل‌ توجه هزینه‌ها نیز می‌شود.

\section{نحوه ارزیابی دستاوردهای تحقیق}
ایجاد یک سامانه دوقلوی مجازی یکپارچه از تعدادی دوقلو‌های مجازی کوچکتر، دستاورد اصلی این پژوهش خواهد بود. در واقع بررسی دوقلو‌های مجازی و ارتباط میان آن‌ها، به طوری که به صورت یکپارچه بتوانند بر روی همتای حقیقی خود نظارت و کنترل انجام دهند، به عنوان یک خروجی از پژوهش شناخته می‌شود.

همچنین شبیه‌سازی و ارزیابی در بستر مدل‌های سیستم‌های دوقلو‌ی مجازی موجود و توسعه‌ یافته، انجام خواهد شد به گونه‌ای که سامانه قابل درک و محسوس باشد.

از آنجایی که هوش‌ مصنوعی جزء جدا نشدنی از دوقلو‌های مجازی است، تلاش در جهت افزایش دقت و سرعت هوش مصنوعی مورد استفاده در سیستم، نسبت به کار‌های مشابه نیز بخشی از اهداف پژوهش می‌باشد.

همچنین تهیه مجموعه داده مورد نیاز و مناسب برای بهبود نتیجه یادگیری دوقلو‌های مجازی در سامانه نیز می‌تواند یکی از دستاورد‌های پژوهش در نظر گرفته شود که تاثیر چشم‌گیری در نتیجه اصلی نیز دارد.



\section{زمان‌بندی تحقیق}
در جدول \ref{tab1}، برنامه‌ی یکسال آینده، زمان‌بندی تقریبی انجام مراحل مختلف تحقیق جهت بدست آوردن نتایج تحقیق برآورد شده است.
\begin{table}[H]
	\centering
	\caption{پیش بینی زمان بندی فعالیت‌ها و مراحل اجرایی تحقیق و ارائه گزارش پیشرفت کار}
	\label{tab1}
	\scalebox{.80}{
		\begin{tabular}{|c|c|c|c|c|c|c|c|c|c|c|c|c|c|c|}
			\hline
			ردیف &
			\multicolumn{1}{>{\centering}m{7cm}|}{شرح فعالیت}
			&	 زمان کل &
			\multicolumn{12}{c|}{زمان اجرا به ماه}\\
			\cline{4-15}
			&
			\multicolumn{1}{>{\centering}m{4cm}|}{}&
			(ماه) & 1 & 2 & 3 & 4 & 5 & 6 & 7 & 8 & 9 & 10 & 11 & 12
			\\ \hline
			1 &
			\multicolumn{1}{>{\centering}m{7cm}|}{مطالعه و پیاده‌سازی سیستم دوقلوی مجازی پایه}	
			& 4 & $\times$ &$\times$ &$\times$ & $\times$& & & & & & & &
			\\ \hline
			2 &
			\multicolumn{1}{>{\centering}m{7cm}|}{ارائه راهکار جدید در ایجاد دوقلوی مجازی}	
			& 5 & & & $\times$ &$\times$ & $\times$&$\times$ &$\times$ & & & & &
			\\ \hline
			3
			&
			\multicolumn{1}{>{\centering}m{7cm}|}{شبیه‌سازی ساختار جدید و ارزیابی نتایج حاصل}
			& 5 & & &  & & & &$\times$ &$\times$ &$\times$ &$\times$ &$\times$ &
			\\ \hline
			4 &
			\multicolumn{1}{>{\centering}m{7cm}|}{نگارش پایان‌نامه}
			& 5 & & & &  &  &  &  & $\times$ & $\times$ & $\times$ & $\times$ & $\times$
			\\ \hline
%			5 &
%			\multicolumn{1}{>{\centering}m{7cm}|}{}
%			& & & & & & & & & & & & &
%			\\ \hline
	\end{tabular}}
\end{table}
\section{امكانات و تجهيزات مورد نياز، نحوه تامين و جدول  برآورد هزينه}
سرور‌های محاسباتی و دارای سخت‌افزار پردازشی جهت پردازش‌های مورد نیاز در یادگیری‌ دوقلو‌های مجازی و هوش‌ مصنوعی.

تعدادی سنسور‌های مختلف مربوط به پارامتر‌های فیزیکی و بورد‌های سخت‌افزار‌ی\LTRfootnote{‌Hardware Boards}
مانند رزبری پای\LTRfootnote{Raspberry Pi}
و آردوینو\LTRfootnote{Arduino}
در جهت شبیه‌سازی دوقلو‌های مجازی.
 
 در صورت امکان، پیاده‌سازی در آزمایشگاه صنعت چهار موجود در دانشگاه زیرا اکثر تجهیزات مورد نیاز در آن وجود دارد.
 
 


\newpage

\renewcommand{\bibname}{مراجع}
\addcontentsline{toc}{chapter}{مراجع}

\bibliographystyle{IEEEtran}

\setLTRbibitems

\begin{thebibliography}{99}

\begin{latin}

\bibitem{b1}
Grieves, M. and Vickers, J. (2017) 
\textit{‘Digital Twin: Mitigating Unpredictable, Undesirable Emergent Behavior in Complex Systems’}
, in Kahlen, J., Flumerfelt, S. and Alves, A. (eds) Transdisciplinary Perspectives on Complex Systems. Cham: Springer. Available at: \url{https://doi.org/10.1007/978-3-319-38756-7_4} 

\bibitem{b2}
Barricelli, B. R., Casiraghi, E. and Fogli, D. (2019) \textit{'A Survey on Digital Twin: Definitions, Characteristics, Applications, and Design Implications'}, IEEE Access, vol. 7, pp. 167653-167671. Available at: \url{https://doi.org/10.1109/ACCESS.2019.2953499}

\bibitem{b3}
Dani, A.A.H., Supangkat, S.H., Lubis, F.F., Nugraha, I.G.B.B., Kinanda, R. and Rizkia, I. (2023) \textit{‘Development of a Smart City Platform Based on Digital Twin Technology for Monitoring and Supporting Decision-Making’}, Sustainability, 15(18), 14002. Available at: \url{https://doi.org/10.3390/su151814002}

\bibitem{b4}
Tan, C., et al. (2023) \textit{‘Digital Twin Enabled Remote Data Sharing for Internet of Vehicles: System and Incentive Design’}, IEEE Transactions on Vehicular Technology, 72(10), pp. 13474-13489. Available at: \url{https://doi.org/10.1109/TVT.2023.3275591}

\bibitem{b5}
Marah, H. and Challenger, M. (2024) \textit{‘MADTwin: a framework for multi-agent digital twin development: smart warehouse case study’}, Annals of Mathematics and Artificial Intelligence, 92, pp. 975–1005. Available at: \url{https://doi.org/10.1007/s10472-023-09872-z}

\bibitem{b6}
Pretel, E., Navarro, E., López-Jaquero, V., Moya, A. and González, P. (2022) \textit{‘Multi-Agent Systems in Support of Digital Twins: A Survey’}, in Ferrández Vicente, J.M., Álvarez-Sánchez, J.R., de la Paz López, F. and Adeli, H. (eds) Bio-inspired Systems and Applications: from Robotics to Ambient Intelligence. IWINAC 2022. Lecture Notes in Computer Science, vol 13259. Springer, Cham, pp.  524–533. Available at: \url{https://doi.org/10.1007/978-3-031-06527-9_52}


\end{latin}

\end{thebibliography}


\newpage
	%%%%%%%% تعهدنامه


\begin{center}

	\scalebox{.80}{
		\begin{tabular}{|c|}
			\hline
			\\
			\multicolumn{1}{|>{\raggedright}m{19cm}|}{\textbf{تعهد دانشجو:}}
			\\
			\\
			%%%%%%%% اطلاعات (فرم را با اطلاعات خود پر کنید.)
			\multicolumn{1}{|>{\raggedright}m{19cm}|}{
				اینجانب ---------------- ضمن اطلاع از اینکه كليه حقوق مادی مترتب بر نتايج تحقيق این پایان نامه متعلق به دانشگاه صنعتی اصفهان بوده و انتشار نتايج تابع مقررات دانشگاهی است که صرفاً پس از موافقت و تایید استاد/ اساتيد راهنما قابل انجام است، تعهد می‌نمایم که اصول اخلاق علمی در انتخاب موضوع و نگارش این گزارش رعایت شده و در طول تحقیق و تدوین/انتشار  نتایج/گزارش ها/مقالات  بعدی نیز رعايت شود.
			}
			\\
			\multicolumn{1}{|>{\raggedright}m{19cm}|}{
				%				نام خود را وارد کنید.
				نام و نام خانوادگي: 
				\hspace{.5cm}
				امضا:
				%			فایل امضای خود را با نام Signature.JPG در مسیر قرار دهید و کامنت بعدی را انکامنت کنید
				%			\includegraphics[scale=.4]{Signature.JPG}
				\hspace{1.5cm}
				تاریخ:
			}
			\\ \hline
	\end{tabular}}
\end{center}
%\clearpage
%%%%%%%%%%%%%%%%%%%%%%%%%%%%%%%%%%%%%%%%%%%%%%%%%%%%%%%%
\begin{center}
	\scalebox{.80}{
		\begin{tabular}{|c|}
			\hline
			\\
			\multicolumn{1}{|>{\raggedright}m{19cm}|}{\textbf{تائيد اساتيد راهنما و مشاور پایان‌نامه:}}
			\\
			\\
		
			%%%%%%%% اطلاعات
			\multicolumn{1}{|>{\raggedright}m{19cm}|}{
				در صورت تصويب موضوع تحقيق پيشنهادي، بدين‌وسيله آمادگي خود را براي راهنمايي و مشاوره دانشجو در كليه مراحل انجام و ارائه تحقيق و مشاركت در ارزيابي پایان‌نامه بر اساس ضوابط دانشگاه و دانشكده اعلام مي‌نماييم. 
			}
			\\
			\multicolumn{1}{|>{\raggedright}m{19cm}|}{
				\hspace{5cm}نام و نام خانوادگی:
				\hspace{2.5cm}
				امضا:
				%			\begin{figure}
					%				\includegraphics[scale=0.4]{Signature.JPG}
					%			\end{figure}
				\hspace{2.5cm}
				تاریخ:
			}
			\\
			
			\multicolumn{1}{|>{\raggedright}m{18cm}|}{
				\begin{enumerate}[1-]
					\item
					استاد راهنما
					\item
					استاد مشاور 
					\item
					استاد راهنمای دوم*
					\item
					استاد مشاور دوم*
				\end{enumerate}
			}
			
			\\ \hline
	\end{tabular}}
\end{center}

%% نظر نهایی گروه

\begin{center}
	\scalebox{.80}{
		\begin{tabular}{|c|}
			\hline
			\\
			\multicolumn{1}{|>{\raggedright}m{19cm}|}{\textbf{نظر نهایی گروه:}}
			\\
			\\
			
			%%%%%%%% اطلاعات
			\multicolumn{1}{|>{\raggedright}m{19cm}|}{
				نسخه نهایی پیشنهاده مورد تایید است.
			}
			\\
			\multicolumn{1}{|>{\raggedright}m{19cm}|}{
				\hspace{5cm}نام و نام خانوادگی مدیر گروه:
				\hspace{2.5cm}
				امضا:
				%			\begin{figure}
					%				\includegraphics[scale=0.4]{Signature.JPG}
					%			\end{figure}
				\hspace{2.5cm}
				تاریخ:
			}
			\\
			
			\\ \hline
	\end{tabular}}
\end{center}

%% نظر نهایی شورای تحصیلات تکمیلی دانشکده

\begin{center}
	\scalebox{.80}{
		\begin{tabular}{|c|}
			\hline
			\\
			\multicolumn{1}{|>{\raggedright}m{19cm}|}{\textbf{نظر نهایی شورای تحصیلات تکمیلی دانشکده:}}
			\\
			\\
			
			%%%%%%%% اطلاعات
			\multicolumn{1}{|>{\raggedright}m{19cm}|}{
موضوع پیشنهادی تحقیق پایان‌نامه کارشناسی ارشد آقا/خانم --------------- مورد تایید است.
			}
			\\
			\multicolumn{1}{|>{\raggedright}m{19cm}|}{
				\hspace{5cm}نام و نام خانوادگی مدیر گروه:
				\hspace{2.5cm}
				امضا:
				%			\begin{figure}
					%				\includegraphics[scale=0.4]{Signature.JPG}
					%			\end{figure}
				\hspace{2.5cm}
				تاریخ:
			}
			\\
			
			\\ \hline
	\end{tabular}}
\end{center}


%\bibliography{REF}
\end{document}
%% End of file 